\section{Conclusiones}

\begin{itemize}
\item En \openbmi{}, la comparación de protocolos y el indicador de degradación mostraron que la diferencia entre personas vistas y no vistas debe distinguirse del cambio de volumen de datos; solo LaBraM mostró degradación significativa tras Holm.
\item Las variantes de procesamiento y alineamiento no tuvieron un efecto común a todas las arquitecturas; la calibración ensayada con pocos ejemplos no produjo ganancia.
\item La transferencia a WBCIC-SHU redujo la exactitud respecto de personas nuevas de \openbmi{}, lo que limita la extrapolación entre protocolos de adquisición.
\item Los criterios adicionales mostraron que la elección de arquitectura depende tanto de la equivalencia práctica como del costo y de las condiciones de medición; la robustez y la saliencia se examinaron en una sola partición y semilla.
\end{itemize}
